% Options for packages loaded elsewhere
\PassOptionsToPackage{unicode}{hyperref}
\PassOptionsToPackage{hyphens}{url}
%
\documentclass[
  ignorenonframetext,
]{beamer}
\usepackage{pgfpages}
\setbeamertemplate{caption}[numbered]
\setbeamertemplate{caption label separator}{: }
\setbeamercolor{caption name}{fg=normal text.fg}
\beamertemplatenavigationsymbolsempty
% Prevent slide breaks in the middle of a paragraph
\widowpenalties 1 10000
\raggedbottom
\setbeamertemplate{part page}{
  \centering
  \begin{beamercolorbox}[sep=16pt,center]{part title}
    \usebeamerfont{part title}\insertpart\par
  \end{beamercolorbox}
}
\setbeamertemplate{section page}{
  \centering
  \begin{beamercolorbox}[sep=12pt,center]{part title}
    \usebeamerfont{section title}\insertsection\par
  \end{beamercolorbox}
}
\setbeamertemplate{subsection page}{
  \centering
  \begin{beamercolorbox}[sep=8pt,center]{part title}
    \usebeamerfont{subsection title}\insertsubsection\par
  \end{beamercolorbox}
}
\AtBeginPart{
  \frame{\partpage}
}
\AtBeginSection{
  \ifbibliography
  \else
    \frame{\sectionpage}
  \fi
}
\AtBeginSubsection{
  \frame{\subsectionpage}
}
\usepackage{amsmath,amssymb}
\usepackage{lmodern}
\usepackage{iftex}
\ifPDFTeX
  \usepackage[T1]{fontenc}
  \usepackage[utf8]{inputenc}
  \usepackage{textcomp} % provide euro and other symbols
\else % if luatex or xetex
  \usepackage{unicode-math}
  \defaultfontfeatures{Scale=MatchLowercase}
  \defaultfontfeatures[\rmfamily]{Ligatures=TeX,Scale=1}
\fi
\usetheme[]{Montpellier}
% Use upquote if available, for straight quotes in verbatim environments
\IfFileExists{upquote.sty}{\usepackage{upquote}}{}
\IfFileExists{microtype.sty}{% use microtype if available
  \usepackage[]{microtype}
  \UseMicrotypeSet[protrusion]{basicmath} % disable protrusion for tt fonts
}{}
\makeatletter
\@ifundefined{KOMAClassName}{% if non-KOMA class
  \IfFileExists{parskip.sty}{%
    \usepackage{parskip}
  }{% else
    \setlength{\parindent}{0pt}
    \setlength{\parskip}{6pt plus 2pt minus 1pt}}
}{% if KOMA class
  \KOMAoptions{parskip=half}}
\makeatother
\usepackage{xcolor}
\IfFileExists{xurl.sty}{\usepackage{xurl}}{} % add URL line breaks if available
\IfFileExists{bookmark.sty}{\usepackage{bookmark}}{\usepackage{hyperref}}
\hypersetup{
  hidelinks,
  pdfcreator={LaTeX via pandoc}}
\urlstyle{same} % disable monospaced font for URLs
\newif\ifbibliography
\usepackage{color}
\usepackage{fancyvrb}
\newcommand{\VerbBar}{|}
\newcommand{\VERB}{\Verb[commandchars=\\\{\}]}
\DefineVerbatimEnvironment{Highlighting}{Verbatim}{commandchars=\\\{\}}
% Add ',fontsize=\small' for more characters per line
\usepackage{framed}
\definecolor{shadecolor}{RGB}{248,248,248}
\newenvironment{Shaded}{\begin{snugshade}}{\end{snugshade}}
\newcommand{\AlertTok}[1]{\textcolor[rgb]{0.94,0.16,0.16}{#1}}
\newcommand{\AnnotationTok}[1]{\textcolor[rgb]{0.56,0.35,0.01}{\textbf{\textit{#1}}}}
\newcommand{\AttributeTok}[1]{\textcolor[rgb]{0.77,0.63,0.00}{#1}}
\newcommand{\BaseNTok}[1]{\textcolor[rgb]{0.00,0.00,0.81}{#1}}
\newcommand{\BuiltInTok}[1]{#1}
\newcommand{\CharTok}[1]{\textcolor[rgb]{0.31,0.60,0.02}{#1}}
\newcommand{\CommentTok}[1]{\textcolor[rgb]{0.56,0.35,0.01}{\textit{#1}}}
\newcommand{\CommentVarTok}[1]{\textcolor[rgb]{0.56,0.35,0.01}{\textbf{\textit{#1}}}}
\newcommand{\ConstantTok}[1]{\textcolor[rgb]{0.00,0.00,0.00}{#1}}
\newcommand{\ControlFlowTok}[1]{\textcolor[rgb]{0.13,0.29,0.53}{\textbf{#1}}}
\newcommand{\DataTypeTok}[1]{\textcolor[rgb]{0.13,0.29,0.53}{#1}}
\newcommand{\DecValTok}[1]{\textcolor[rgb]{0.00,0.00,0.81}{#1}}
\newcommand{\DocumentationTok}[1]{\textcolor[rgb]{0.56,0.35,0.01}{\textbf{\textit{#1}}}}
\newcommand{\ErrorTok}[1]{\textcolor[rgb]{0.64,0.00,0.00}{\textbf{#1}}}
\newcommand{\ExtensionTok}[1]{#1}
\newcommand{\FloatTok}[1]{\textcolor[rgb]{0.00,0.00,0.81}{#1}}
\newcommand{\FunctionTok}[1]{\textcolor[rgb]{0.00,0.00,0.00}{#1}}
\newcommand{\ImportTok}[1]{#1}
\newcommand{\InformationTok}[1]{\textcolor[rgb]{0.56,0.35,0.01}{\textbf{\textit{#1}}}}
\newcommand{\KeywordTok}[1]{\textcolor[rgb]{0.13,0.29,0.53}{\textbf{#1}}}
\newcommand{\NormalTok}[1]{#1}
\newcommand{\OperatorTok}[1]{\textcolor[rgb]{0.81,0.36,0.00}{\textbf{#1}}}
\newcommand{\OtherTok}[1]{\textcolor[rgb]{0.56,0.35,0.01}{#1}}
\newcommand{\PreprocessorTok}[1]{\textcolor[rgb]{0.56,0.35,0.01}{\textit{#1}}}
\newcommand{\RegionMarkerTok}[1]{#1}
\newcommand{\SpecialCharTok}[1]{\textcolor[rgb]{0.00,0.00,0.00}{#1}}
\newcommand{\SpecialStringTok}[1]{\textcolor[rgb]{0.31,0.60,0.02}{#1}}
\newcommand{\StringTok}[1]{\textcolor[rgb]{0.31,0.60,0.02}{#1}}
\newcommand{\VariableTok}[1]{\textcolor[rgb]{0.00,0.00,0.00}{#1}}
\newcommand{\VerbatimStringTok}[1]{\textcolor[rgb]{0.31,0.60,0.02}{#1}}
\newcommand{\WarningTok}[1]{\textcolor[rgb]{0.56,0.35,0.01}{\textbf{\textit{#1}}}}
\setlength{\emergencystretch}{3em} % prevent overfull lines
\providecommand{\tightlist}{%
  \setlength{\itemsep}{0pt}\setlength{\parskip}{0pt}}
\setcounter{secnumdepth}{-\maxdimen} % remove section numbering
%\documentclass[unicode,10pt]{beamer}% 'unicode'が必要
\usepackage{luatexja}% 日本語を使う場合は必須
%\usepackage[japanese]{babel}	
%\usefonttheme[onlymath]{serif}
%\usepackage[T1]{fontenc} %これを使うと、ギリシャ文字が出ない
%\usepackage{textcomp}
%\usepackage[scale=1.0]{tgheros} % Sans serif font
%\usepackage[scaled]{beramono} % Monospace font
%\usepackage{luatexja-otf}
%\renewcommand{\kanjifamilydefault}{\gtdefault}
%\usepackage[haranoaji]{luatexja-preset}

% スライド番号の表示 %%%%%%%%%%%%%%%%%%%
\makeatletter
\setbeamertemplate{footline}{%
  \leavevmode%
  \hbox{%
  \begin{beamercolorbox}[wd=.5\paperwidth,ht=2.25ex,dp=1ex,right]{date in head/foot}%
    \insertframenumber{} \hspace*{2ex} % new version without total frames
  \end{beamercolorbox}}%
  \vskip0pt%
}
\makeatother
% スライド番号の表示 %%%%%%%%%%%%%%%%%%%
\ifLuaTeX
  \usepackage{selnolig}  % disable illegal ligatures
\fi

\title{第3章: 測定\\
(3.2. Rで欠損データを扱う)}
\subtitle{今井耕介 著\\
『社会科学のためのデータ分析入門 (QSS)』}
\author{}
\date{\vspace{-2.5em}2026-03-09}

\begin{document}
\frame{\titlepage}

\hypertarget{rux3067ux6b20ux640dux30c7ux30fcux30bfux3092ux6271ux3046}{%
\section{3.2
Rで欠損データを扱う}\label{rux3067ux6b20ux640dux30c7ux30fcux30bfux3092ux6271ux3046}}

\begin{frame}[fragile]{欠損値 (Missing Data) とは}
\protect\hypertarget{ux6b20ux640dux5024-missing-data-ux3068ux306f}{}
\begin{itemize}
\tightlist
\item
  調査において、回答者が質問に答えなかったり、回答を拒否したりした場合に生じる。
\item
  Rでは、欠損値は \textbf{\texttt{NA}} (Not Available)
  という特別な記号で表される。
\item
  \textbf{注意}: \texttt{NA}
  が含まれるデータに対してそのまま計算（平均など）を行うと、結果も
  \texttt{NA} になってしまう。
\end{itemize}
\end{frame}

\hypertarget{ux6b20ux640dux5024ux306eux78baux8a8d}{%
\section{3.2.1
欠損値の確認}\label{ux6b20ux640dux5024ux306eux78baux8a8d}}

\begin{frame}[fragile]{is.na() 関数}
\protect\hypertarget{is.na-ux95a2ux6570}{}
\begin{itemize}
\tightlist
\item
  データが \texttt{NA} かどうかを判定します。
\end{itemize}

\scriptsize

\begin{Shaded}
\begin{Highlighting}[]
\CommentTok{\# アフガニスタン調査データの読み込み}
\NormalTok{afghan }\OtherTok{\textless{}{-}} \FunctionTok{read.csv}\NormalTok{(}\StringTok{"afghan.csv"}\NormalTok{)}

\CommentTok{\# 1. 最初の10人分の収入データを確認}
\FunctionTok{head}\NormalTok{(afghan}\SpecialCharTok{$}\NormalTok{income, }\AttributeTok{n =} \DecValTok{10}\NormalTok{)}
\end{Highlighting}
\end{Shaded}

\begin{verbatim}
##  [1] "2,001-10,000"  "2,001-10,000"  "2,001-10,000"  "2,001-10,000" 
##  [5] "2,001-10,000"  NA              "10,001-20,000" "2,001-10,000" 
##  [9] "2,001-10,000"  NA
\end{verbatim}

\begin{Shaded}
\begin{Highlighting}[]
\CommentTok{\# 2. 欠損値かどうかを判定 (TRUEが欠損)}
\FunctionTok{head}\NormalTok{(}\FunctionTok{is.na}\NormalTok{(afghan}\SpecialCharTok{$}\NormalTok{income), }\AttributeTok{n =} \DecValTok{10}\NormalTok{)}
\end{Highlighting}
\end{Shaded}

\begin{verbatim}
##  [1] FALSE FALSE FALSE FALSE FALSE  TRUE FALSE FALSE FALSE  TRUE
\end{verbatim}

\normalsize
\end{frame}

\begin{frame}[fragile]{欠損値の数と割合の計算}
\protect\hypertarget{ux6b20ux640dux5024ux306eux6570ux3068ux5272ux5408ux306eux8a08ux7b97}{}
\begin{itemize}
\tightlist
\item
  \texttt{is.na()} の結果（TRUE/FALSE）を \texttt{sum()} や
  \texttt{mean()} に渡すことで集計できます。
\end{itemize}

\scriptsize

\begin{Shaded}
\begin{Highlighting}[]
\CommentTok{\# 1. 欠損値の総数をカウント}
\FunctionTok{sum}\NormalTok{(}\FunctionTok{is.na}\NormalTok{(afghan}\SpecialCharTok{$}\NormalTok{income))}
\end{Highlighting}
\end{Shaded}

\begin{verbatim}
## [1] 154
\end{verbatim}

\begin{Shaded}
\begin{Highlighting}[]
\CommentTok{\# 2. 全体に占める欠損値の割合を計算}
\FunctionTok{mean}\NormalTok{(}\FunctionTok{is.na}\NormalTok{(afghan}\SpecialCharTok{$}\NormalTok{income))}
\end{Highlighting}
\end{Shaded}

\begin{verbatim}
## [1] 0.05591866
\end{verbatim}

\normalsize
\end{frame}

\hypertarget{ux6b20ux640dux5024ux304cux3042ux308bux5834ux5408ux306eux8a08ux7b97}{%
\section{3.2.2
欠損値がある場合の計算}\label{ux6b20ux640dux5024ux304cux3042ux308bux5834ux5408ux306eux8a08ux7b97}}

\begin{frame}[fragile]{na.rm オプション}
\protect\hypertarget{na.rm-ux30aaux30d7ux30b7ux30e7ux30f3}{}
\begin{itemize}
\tightlist
\item
  多くの統計関数には、\texttt{NA} を除外して計算するための
  \texttt{na.rm\ =\ TRUE} という引数があります。
\end{itemize}

\scriptsize

\begin{Shaded}
\begin{Highlighting}[]
\CommentTok{\# サンプルベクトルの作成}
\NormalTok{x }\OtherTok{\textless{}{-}} \FunctionTok{c}\NormalTok{(}\DecValTok{1}\NormalTok{, }\DecValTok{2}\NormalTok{, }\DecValTok{3}\NormalTok{, }\ConstantTok{NA}\NormalTok{)}

\CommentTok{\# そのまま平均を計算すると NA になる}
\FunctionTok{mean}\NormalTok{(x)}
\end{Highlighting}
\end{Shaded}

\begin{verbatim}
## [1] NA
\end{verbatim}

\begin{Shaded}
\begin{Highlighting}[]
\CommentTok{\# NAを除外して平均を計算}
\FunctionTok{mean}\NormalTok{(x, }\AttributeTok{na.rm =} \ConstantTok{TRUE}\NormalTok{)}
\end{Highlighting}
\end{Shaded}

\begin{verbatim}
## [1] 2
\end{verbatim}

\normalsize
\end{frame}

\begin{frame}[fragile]{集計表における欠損値の表示}
\protect\hypertarget{ux96c6ux8a08ux8868ux306bux304aux3051ux308bux6b20ux640dux5024ux306eux8868ux793a}{}
\begin{itemize}
\tightlist
\item
  \texttt{table()} 関数で \texttt{exclude\ =\ NULL}
  を指定すると、\texttt{NA} も一つのカテゴリとして表示されます。
\end{itemize}

\scriptsize

\begin{Shaded}
\begin{Highlighting}[]
\CommentTok{\# 被害経験のクロス集計に NA も含める}
\CommentTok{\# exclude = NULL を入れないと NA は無視される}
\FunctionTok{table}\NormalTok{(}\AttributeTok{ISAF =}\NormalTok{ afghan}\SpecialCharTok{$}\NormalTok{violent.exp.ISAF, }
      \AttributeTok{Taliban =}\NormalTok{ afghan}\SpecialCharTok{$}\NormalTok{violent.exp.taliban, }
      \AttributeTok{exclude =} \ConstantTok{NULL}\NormalTok{)}
\end{Highlighting}
\end{Shaded}

\begin{verbatim}
##       Taliban
## ISAF      0    1 <NA>
##   0    1330  354   22
##   1     475  526   22
##   <NA>    7    8   10
\end{verbatim}

\normalsize
\end{frame}

\hypertarget{ux6b20ux640dux5024ux306eux9664ux5916}{%
\section{3.2.3
欠損値の除外}\label{ux6b20ux640dux5024ux306eux9664ux5916}}

\begin{frame}[fragile]{na.omit() による一括除外}
\protect\hypertarget{na.omit-ux306bux3088ux308bux4e00ux62ecux9664ux5916}{}
\begin{itemize}
\tightlist
\item
  \textbf{リストワイズ削除}: 1つでも \texttt{NA}
  を含む行（回答者）をデータセット全体から削除します。
\end{itemize}

\scriptsize

\begin{Shaded}
\begin{Highlighting}[]
\CommentTok{\# 1. 1つでも欠損値がある行をすべて削除した新しいデータセットを作成}
\NormalTok{afghan.sub }\OtherTok{\textless{}{-}} \FunctionTok{na.omit}\NormalTok{(afghan)}

\CommentTok{\# 2. 元の行数と、削除後の行数を比較}
\FunctionTok{nrow}\NormalTok{(afghan)     }\CommentTok{\# 元のデータ}
\end{Highlighting}
\end{Shaded}

\begin{verbatim}
## [1] 2754
\end{verbatim}

\begin{Shaded}
\begin{Highlighting}[]
\FunctionTok{nrow}\NormalTok{(afghan.sub) }\CommentTok{\# 削除後のデータ}
\end{Highlighting}
\end{Shaded}

\begin{verbatim}
## [1] 2554
\end{verbatim}

\begin{Shaded}
\begin{Highlighting}[]
\CommentTok{\# 3. 特定の変数だけから NA を除く}
\FunctionTok{length}\NormalTok{(}\FunctionTok{na.omit}\NormalTok{(afghan}\SpecialCharTok{$}\NormalTok{income))}
\end{Highlighting}
\end{Shaded}

\begin{verbatim}
## [1] 2600
\end{verbatim}

\normalsize
\end{frame}

\hypertarget{ux307eux3068ux3081}{%
\section{3.2.4 まとめ}\label{ux307eux3068ux3081}}

\begin{frame}[fragile]{このセクションのまとめ}
\protect\hypertarget{ux3053ux306eux30bbux30afux30b7ux30e7ux30f3ux306eux307eux3068ux3081}{}
\begin{itemize}
\tightlist
\item
  \textbf{NAの理解}: Rにおいて欠損値は \texttt{NA}
  で表され、計算に影響を及ぼす。
\item
  \textbf{確認と集計}:

  \begin{itemize}
  \tightlist
  \item
    \texttt{is.na()} で確認し、\texttt{sum()} や \texttt{mean()}
    で量を確認する。
  \item
    \texttt{table(...,\ exclude\ =\ NULL)}
    で回答拒否なども含めた全体像を把握する。
  \end{itemize}
\item
  \textbf{処理方法}:

  \begin{itemize}
  \tightlist
  \item
    \texttt{mean(...,\ na.rm\ =\ TRUE)} のように、その都度除外する。
  \item
    \texttt{na.omit()}
    でデータセットから欠損のある行を一括削除する（分析対象が減ることに注意）。
  \end{itemize}
\item
  \textbf{Rの操作}: \texttt{is.na()}, \texttt{na.rm},
  \texttt{na.omit()}, \texttt{nrow()}。
\end{itemize}
\end{frame}

\end{document}
